\documentclass[11pt,a4paper]{article}
\usepackage[T1]{fontenc}
\usepackage{lmodern}
\usepackage[margin=27mm]{geometry}
\usepackage{amsmath,amssymb,amsthm,mathtools}
\usepackage{microtype}
\usepackage{booktabs,array}
\usepackage{enumitem}
\usepackage[hidelinks]{hyperref}
\setlength{\parskip}{5pt plus 1pt}
\setlength{\parindent}{0pt}
\setlength{\emergencystretch}{3em}
\setlist{itemsep=2pt,topsep=3pt}
\raggedbottom
\newtheorem{theorem}{Theorem}[section]
\newtheorem{lemma}[theorem]{Lemma}
\newtheorem{corollary}[theorem]{Corollary}
\newtheorem{proposition}[theorem]{Proposition}
\theoremstyle{definition}
\newtheorem{definition}[theorem]{Definition}
\newtheorem{example}[theorem]{Example}
\theoremstyle{remark}
\newtheorem{remark}[theorem]{Remark}
\newcommand{\Z}{\mathbb{Z}}
\newcommand{\N}{\mathbb{N}}
\newcommand{\Q}{\mathbb{Q}}
\newcommand{\F}{\mathbb{F}}
\newcommand{\ps}[1]{#1[[x]]}
\newcommand{\coeff}[2]{[x^{#1}]\,#2}
\newcommand{\hl}[1]{\textbf{#1}}
\newcommand{\para}[1]{\par\medskip\noindent\textbf{#1}\enspace}
\hypersetup{pdftitle={The coefficients of (exp(x)-1)/cos(2x) modulo 4 and 5: a proof of two conjectures in OEIS A380056},pdfauthor={Dev D. Goyal},pdfsubject={OEIS A380056}}
\title{The coefficients of $(e^x-1)/\cos 2x$ modulo 4 and 5\\[6pt]\large A proof of two conjectures in OEIS A380056}
\author{Dev D. Goyal}
\date{October 2026}
\begin{document}
\maketitle
\begin{abstract}
The OEIS entry A380056, contributed by Paul D.~Hanna, lists the numbers $a(n)$ defined by the exponential generating
function $(e^x-1)/\cos 2x$; that is, $a(n)/n!$ is the coefficient of $x^n$ in the power series of this function. The
first terms are $1,1,13,25,441,1261,\dots$. Hanna conjectured that $a(n)$ leaves remainder $1$ on division by $4$ for
every $n\ge1$, and that $a(4n)$ is divisible by $5$. We prove both, and more: modulo $5$,
$a(n)\equiv2\cdot2^n+3\cdot4^n$, so the remainders repeat in the cycle $1,1,3,0$, and $a(n)$ is divisible by $5$
exactly when $n$ is a multiple of $4$. The proof first rewrites the definition as a recurrence that uses only whole
numbers, and then solves this recurrence modulo $4$ and modulo $5$. Modulo $5$ the solution is a combination of two
geometric sequences, because there $\cos 2x$ behaves like $(e^x+e^{-x})/2$. The results have also been checked by
computer and proved formally in the Lean proof assistant.
\end{abstract}
\begin{center}\begin{minipage}{0.88\linewidth}\small
\textbf{How this paper was produced.} This paper reports the result of an autonomous AI research effort.
ChatGPT (OpenAI), using the GPT-6 model with Pro reasoning and working from prompts by Dev~D.~Goyal, found the proof, wrote its Lean~4 formalization, and wrote a first draft of the paper.
Dev~D.~Goyal chose the problems and asked for checks and revisions, but did not supply proof ideas.
Claude Opus~5.5, used with Claude Code, then checked the work independently: it rebuilt the formal proofs from
source and confirmed that they rely on no unproved assumptions, recomputed the sequence from its definition and tested
every claim against the published data, and read every step of the proofs. It also rewrote this paper for clarity.
Dev~D.~Goyal has verified the results and is responsible for them. The paper has not been peer reviewed.
The section ``About this paper'' at the end gives the details.
\end{minipage}\end{center}

\section{Introduction}

\subsection{The sequence}\label{sec:sequence}
The \emph{On-Line Encyclopedia of Integer Sequences} (OEIS)~\cite{oeisdb} is a large public database of integer
sequences, used by mathematicians and scientists to identify sequences and to record what is known about them. Many
entries also record \emph{conjectures}: patterns that the contributor observed in the data but could not prove.

This paper is about entry A380056~\cite{oeis}, contributed by Paul~D.~Hanna in January 2025. Its name is:
\begin{quote}\small\ttfamily\raggedright E.g.f. (exp(x) - 1)/cos(2*x).\par\smallskip\hfill\textrm{--- OEIS A380056, Paul D. Hanna, January 2025}\end{quote}

``E.g.f.''\ stands for \emph{exponential generating function} (Section~\ref{sec:background}). The name means that
\begin{equation}\label{eq:def}
\sum_{n\ge0}a(n)\,\frac{x^n}{n!}=\frac{e^x-1}{\cos 2x},
\end{equation}
where $e^x=1+x+\frac{x^2}{2!}+\cdots$ and $\cos2x=1-\frac{(2x)^2}{2!}+\frac{(2x)^4}{4!}-\cdots$ are the usual power series.
In other words, expand $(e^x-1)/\cos2x$ as a power series; the coefficient of $x^n$, multiplied by $n!$, is $a(n)$. The
constant term is $0$, so $a(0)=0$; the entry starts at $n=1$. The first terms, from $n=1$, are
1, 1, 13, 25, 441, 1261, 30213, 115025, 3529201, 16792021, \dots{} It is not obvious from the definition that these numbers are
integers; Lemma~\ref{lem:rec} below shows that they are.

\subsection{The question}
The entry contains two conjectures.
\begin{quote}\small\ttfamily\raggedright Conjecture: a(n) == 1 (mod 4) for n >= 1.\par\smallskip Conjecture: a(4*n) is divisible by 5 for n >= 1.\par\smallskip\hfill\textrm{--- OEIS A380056, Paul D. Hanna, January 2025}\end{quote}

Here \texttt{a(n) == 1 (mod 4)} means that $a(n)$ leaves remainder $1$ on division by $4$. The table shows the
remainders of the first terms modulo $4$ and modulo $5$. The indices in bold are the multiples of $4$.
\begin{center}\small\setlength{\tabcolsep}{3.2pt}
\begin{tabular}{>{\raggedright\arraybackslash}p{6.4em}>{\centering\arraybackslash}p{1.15em}>{\centering\arraybackslash}p{1.15em}>{\centering\arraybackslash}p{1.15em}>{\centering\arraybackslash}p{1.15em}>{\centering\arraybackslash}p{1.15em}>{\centering\arraybackslash}p{1.15em}>{\centering\arraybackslash}p{1.15em}>{\centering\arraybackslash}p{1.15em}>{\centering\arraybackslash}p{1.15em}>{\centering\arraybackslash}p{1.15em}>{\centering\arraybackslash}p{1.15em}>{\centering\arraybackslash}p{1.15em}>{\centering\arraybackslash}p{1.15em}>{\centering\arraybackslash}p{1.15em}>{\centering\arraybackslash}p{1.15em}>{\centering\arraybackslash}p{1.15em}}
\toprule
$n$ & 1 & 2 & 3 & \hl{4} & 5 & 6 & 7 & \hl{8} & 9 & 10 & 11 & \hl{12} & 13 & 14 & 15 & \hl{16} \\ \midrule
$a(n)\bmod 4$ & 1 & 1 & 1 & \hl{1} & 1 & 1 & 1 & \hl{1} & 1 & 1 & 1 & \hl{1} & 1 & 1 & 1 & \hl{1} \\
$a(n)\bmod 5$ & 1 & 1 & 3 & \hl{0} & 1 & 1 & 3 & \hl{0} & 1 & 1 & 3 & \hl{0} & 1 & 1 & 3 & \hl{0} \\
\bottomrule
\end{tabular}
\\[4pt]
\begin{tabular}{>{\raggedright\arraybackslash}p{6.4em}>{\centering\arraybackslash}p{1.15em}>{\centering\arraybackslash}p{1.15em}>{\centering\arraybackslash}p{1.15em}>{\centering\arraybackslash}p{1.15em}>{\centering\arraybackslash}p{1.15em}>{\centering\arraybackslash}p{1.15em}>{\centering\arraybackslash}p{1.15em}>{\centering\arraybackslash}p{1.15em}>{\centering\arraybackslash}p{1.15em}>{\centering\arraybackslash}p{1.15em}>{\centering\arraybackslash}p{1.15em}>{\centering\arraybackslash}p{1.15em}>{\centering\arraybackslash}p{1.15em}>{\centering\arraybackslash}p{1.15em}>{\centering\arraybackslash}p{1.15em}>{\centering\arraybackslash}p{1.15em}}
\toprule
$n$ & 17 & 18 & 19 & \hl{20} & 21 & 22 & 23 & \hl{24} & 25 & 26 & 27 & \hl{28} & 29 & 30 & 31 & \hl{32} \\ \midrule
$a(n)\bmod 4$ & 1 & 1 & 1 & \hl{1} & 1 & 1 & 1 & \hl{1} & 1 & 1 & 1 & \hl{1} & 1 & 1 & 1 & \hl{1} \\
$a(n)\bmod 5$ & 1 & 1 & 3 & \hl{0} & 1 & 1 & 3 & \hl{0} & 1 & 1 & 3 & \hl{0} & 1 & 1 & 3 & \hl{0} \\
\bottomrule
\end{tabular}
\end{center}

Modulo $5$ the data show more than the conjecture: the remainders repeat in the cycle $1,1,3,0$.

\subsection{What this paper proves}
\textbf{Main result} (Theorem~\ref{thm:main}). \emph{The numbers $a(n)$ are integers, and}
\begin{itemize}
\item \emph{$a(n)\equiv1\pmod4$ for every $n\ge1$;}
\item \emph{$a(n)\equiv2\cdot2^n+3\cdot4^n\pmod5$ for every $n\ge0$. So the remainder of $a(n)$ modulo $5$ is $0,1,1,3$
according as $n$ leaves remainder $0,1,2,3$ on division by $4$; in particular $a(n)$ is divisible by $5$ exactly when
$n$ is a multiple of $4$.}
\end{itemize}
Both conjectures follow (Corollary~\ref{cor:conj}).

\subsection{What was known}
The entry records, among other formulas, an expression for $a(n)$ as a sum involving the Euler (secant) numbers
$1,1,5,61,1385,\dots$ (OEIS A000364~\cite{oeissec}) and powers of $4$ (its formulas~(3.a) and~(3.b)), and an asymptotic
formula due to Vaclav Kotesovec. That expression shows that the $a(n)$ are integers. It also implies the first
conjecture directly, because every term of the sum except one contains a factor $4$, and the remaining term is $1$.
(We checked this numerically as well.) So the first conjecture was, in effect, within reach of the information in the
entry. We are not aware of an earlier proof of the second conjecture, or of the formula modulo $5$.

\subsection{Why it is interesting}
Exponential generating functions are the natural way to encode sequences that count arrangements of labelled objects,
and they appear throughout combinatorics. Their coefficients $n!\cdot\coeff{n}{F}$ are often integers, even though the
series $F$ itself has fractional coefficients, and the remainders of such integers modulo small numbers often follow
patterns. Congruences of this kind for the closely related Euler numbers go back to the nineteenth century. The
difficulty is always the same: the factorials in the denominators prevent us from simply reducing everything modulo
$m$. This paper shows a clean way around it (Lemma~\ref{lem:conv}), and an example where, modulo $5$, a trigonometric
function turns into a combination of exponentials.

\subsection{How to read this paper}
Section~\ref{sec:background} explains the notions we need: power series, remainders, exponential generating
functions, and a way to multiply them without fractions. Section~\ref{sec:results} states the results precisely, and
Section~\ref{sec:proofs} proves them. Section~\ref{sec:checks} describes how the results were checked, by computer
and by a formal proof in Lean. Section~\ref{sec:discussion} discusses other moduli.
\section{Background}\label{sec:background}
This section collects the definitions and facts used later. Readers who know formal power series and congruences can skip it and return when needed.
\subsection{Formal power series}
A \emph{formal power series} is an expression
\[
U(x)=u_0+u_1x+u_2x^2+u_3x^3+\cdots ,
\]
determined by its sequence of coefficients $u_0,u_1,u_2,\dots$. No number is ever substituted for $x$, so questions
of convergence do not arise. We write $\coeff{n}{U}=u_n$ for the coefficient of $x^n$ in $U$; two series are equal
when all their coefficients are equal.

Series are added coefficient by coefficient and multiplied like polynomials:
\[
\coeff{n}{(UV)}=u_0v_n+u_1v_{n-1}+\cdots+u_nv_0 .
\]
Each coefficient of a product is a finite sum. For example, $(1-x)(1+x+x^2+x^3+\cdots)=1$, so $1+x+x^2+\cdots$ is
the \emph{multiplicative inverse} of $1-x$, written $\frac{1}{1-x}$. More generally, every series $D$ with constant term
$1$ has a unique multiplicative inverse $\frac1D$, whose coefficients are determined one at a time by $D\cdot\frac1D=1$
using only addition, subtraction and multiplication. Every fraction in this paper has a denominator of this kind, so it
denotes a power series.

We write $U=O(x^d)$ when all coefficients of $U$ below degree $d$ are zero. So
$U-V=O(x^d)$ means that $U$ and $V$ agree up to degree $d-1$. A series of the form $x+u_2x^2+u_3x^3+\cdots$, with
constant term $0$ and coefficient of $x$ equal to $1$, is called \emph{normalized}.

\subsection{Remainders and congruences}
For an integer $m\ge2$ we write $a\equiv b\pmod m$ (``$a$ is congruent to $b$ modulo $m$'') when $a-b$ is divisible by
$m$, that is, when $a$ and $b$ leave the same remainder after division by $m$. For example, $17\equiv5\pmod 4$, since both leave remainder $1$, and $-3\equiv1\pmod4$. The
remainder of $a$ on division by $m$ is called its \emph{residue} modulo $m$; it is one of $0,1,\dots,m-1$. Saying that
$a$ is odd is the same as saying $a\equiv1\pmod2$.

Addition and multiplication are compatible with congruences, so the residues $0,1,\dots,m-1$, added and multiplied
modulo $m$, form a ring, written $\Z/m\Z$. When $m=p$ is a prime we also write $\F_p$. In $\F_2=\{0,1\}$ we have $1+1=0$,
so addition and subtraction coincide.

For power series with integer coefficients, $U\equiv V\pmod m$ means that \emph{every} coefficient of $U-V$ is
divisible by $m$. Equivalently, the two series become equal when every coefficient is replaced by its residue; we call
this \emph{reducing} the series modulo $m$. The key fact is the following.

\emph{Reduction modulo $m$ is compatible with all the operations above.} To reduce a sum, a product, a substitution, or
an inverse, we may reduce the ingredients first and then compute in $\Z/m\Z$. This holds because every
coefficient of the result is obtained from finitely many coefficients of the ingredients by adding, subtracting and
multiplying, and these operations are compatible with remainders.

\para{Rings.} Several lemmas below hold over every \emph{commutative ring}, that is, every set with an addition,
subtraction and multiplication satisfying the usual rules of arithmetic (for example $\Z$, $\Q$ and $\Z/m\Z$). No
division is used; this matters because in $\Z/m\Z$ division is not always possible (for instance, $2$ has no inverse in
$\Z/4\Z$).

\subsection{Exponential generating functions}
The \emph{exponential generating function} (e.g.f.) of a sequence $f(0),f(1),f(2),\dots$ is the power series
\[
\mathcal E(f)=\sum_{n\ge0}f(n)\,\frac{x^n}{n!}.
\]
The division by $n!$ is a convention that suits sequences which count arrangements of $n$ labelled objects; such
sequences grow roughly like $n!$, and the division keeps the coefficients manageable. Two examples matter here. The
constant sequence $1,1,1,\dots$ has e.g.f.\ $\sum_nx^n/n!=e^x$, the exponential series. And the sequence $1,0,-1,0,1,0,\dots$
has e.g.f.\ $1-\frac{x^2}{2!}+\frac{x^4}{4!}-\cdots=\cos x$. We use these only as formal power series: $e^x$ and
$\cos x$ are defined by these series, and $x$ is never given a value.

Replacing $x$ by $2x$ multiplies the coefficient of $x^n$ by $2^n$. So
\[
\cos2x=\sum_{j\ge0}(-1)^j\frac{(2x)^{2j}}{(2j)!}=\sum_{j\ge0}(-4)^j\,\frac{x^{2j}}{(2j)!},
\]
that is, $\cos 2x$ is the e.g.f.\ of the sequence $c$ with $c(2j)=(-4)^j$ and $c(2j+1)=0$:
$c=1,0,-4,0,16,0,-64,\dots$. Its constant term is $1$, so $\cos2x$ has a multiplicative inverse among the power series
with rational coefficients, and the quotient $(e^x-1)/\cos2x$ is a power series.

\subsection{Multiplying e.g.f.s without fractions}
\begin{lemma}[Binomial convolution]\label{lem:conv}
For two sequences $f$ and $g$, define the sequence $f\star g$ by
\[
(f\star g)(n)=\sum_{k=0}^{n}\binom nk\,f(k)\,g(n-k).
\]
Then $\mathcal E(f)\cdot\mathcal E(g)=\mathcal E(f\star g)$.
\end{lemma}

\begin{proof}
The coefficient of $x^n$ in the product is $\sum_k\frac{f(k)}{k!}\cdot\frac{g(n-k)}{(n-k)!}$. Multiply by $n!$ and use
$\binom nk=\frac{n!}{k!\,(n-k)!}$.
\end{proof}

The point of this lemma is that $f\star g$ is computed from $f$ and $g$ with additions and multiplications of whole
numbers only. So it makes sense for sequences of residues modulo any $m$, and reduction modulo $m$ is compatible with
it. The e.g.f.\ itself, with its denominators $n!$, cannot be reduced modulo $m$ directly.

Two simple facts about $\star$ will be used. First, the sequence $\delta=1,0,0,0,\dots$ is neutral:
$\delta\star f=f$ for every $f$. Second, write $E_s$ for the geometric sequence $E_s(n)=s^n$ (with $s^0=1$). Then
\[
E_s\star E_t=E_{s+t},
\]
because $\sum_k\binom nks^kt^{n-k}=(s+t)^n$ by the binomial theorem. (In e.g.f.\ language this is $e^{sx}e^{tx}=e^{(s+t)x}$.)
These facts hold for numbers $s,t$ in any commutative ring, in particular for residues modulo $5$.

\section{Results}\label{sec:results}

\begin{theorem}\label{thm:main}
The numbers $a(n)$ defined by~\eqref{eq:def} are integers, and:
\begin{enumerate}[label=(\alph*)]
\item $a(n)\equiv1\pmod4$ for every $n\ge1$;
\item $a(n)\equiv2\cdot2^n+3\cdot4^n\pmod5$ for every $n\ge0$. Equivalently, $a(n)\bmod5$ is $0,1,1,3$ according as
$n\equiv0,1,2,3\pmod4$. In particular, $5$ divides $a(n)$ if and only if $4$ divides $n$.
\end{enumerate}
\end{theorem}

\begin{corollary}[Hanna's conjectures]\label{cor:conj}
$a(n)\equiv1\pmod 4$ for $n\ge1$, and $a(4n)$ is divisible by $5$ for every $n\ge1$.
\end{corollary}

\begin{proof}
The first statement is part~(a). The second is the case ``$4$ divides the index'' of part~(b).
\end{proof}

\section{Proofs}\label{sec:proofs}

\para{The idea in brief.} By Lemma~\ref{lem:conv}, multiplying~\eqref{eq:def} by $\cos2x$ turns the definition into a
recurrence for $a(n)$ that uses only whole numbers. Modulo $4$, every coefficient of $\cos 2x$ except the first is
divisible by $4$, so $\cos2x$ acts like $1$ and $a$ is simply the e.g.f.\ sequence of $e^x-1$. Modulo $5$ we have
$-4\equiv1$, so the coefficients of $\cos2x$ become those of $(e^x+e^{-x})/2$. Everything is then a combination of
exponentials $e^{sx}$, whose coefficient sequences are the geometric sequences $s^n$, and a short calculation finds
the answer.

Throughout, let $c$ be the coefficient sequence of $\cos 2x$ (Section~\ref{sec:background}), and let
$d=0,1,1,1,\dots$ be that of $e^x-1$, so that $\mathcal E(c)=\cos 2x$ and $\mathcal E(d)=e^x-1$.

\begin{lemma}\label{lem:rec}
In every commutative ring there is exactly one sequence $a$ with $a\star c=d$, namely the one given by $a(0)=0$ and
\[
a(n)=1-\sum_{k=0}^{n-1}\binom nk\,a(k)\,c(n-k)\qquad(n\ge1).
\]
Over the integers this is the sequence of OEIS A380056. In particular its terms are integers, and reducing it
modulo $m$ gives the solution of $a\star c=d$ over $\Z/m\Z$.
\end{lemma}

\begin{proof}
Since $c(0)=1$, the equation $(a\star c)(n)=d(n)$ reads $a(n)+\sum_{k<n}\binom nk a(k)c(n-k)=d(n)$, which is the
displayed recurrence (with $d(0)=0$ and $d(n)=1$ for $n\ge1$). It determines each $a(n)$ from the earlier terms by
additions, subtractions and multiplications, so it has exactly one solution in every commutative ring, and reduction
modulo $m$ maps one solution to the other. Over the rational numbers, Lemma~\ref{lem:conv} gives
$\mathcal E(a)\cdot\cos2x=\mathcal E(a\star c)=e^x-1$, and dividing by $\cos2x$ shows that $\mathcal E(a)$ is the series
in~\eqref{eq:def}. So the integer solution is the sequence of the entry.
\end{proof}

\begin{proof}[Proof of Theorem~\ref{thm:main}]
The integrality was shown in Lemma~\ref{lem:rec}.

(a) Modulo $4$, every term $c(2j)=(-4)^j$ with $j\ge1$ is $0$, and the odd terms are $0$ anyway. So $c\equiv\delta\pmod4$,
the neutral sequence $1,0,0,\dots$. The sequence $d$ solves $d\star\delta=d$, so by the uniqueness in
Lemma~\ref{lem:rec}, applied over $\Z/4\Z$, we get $a\equiv d\pmod4$. That is, $a(n)\equiv1\pmod4$ for every $n\ge1$.

(b) Work in $\F_5$, the residues modulo $5$. There $-4=1$, so $c(n)=1$ for even $n$ and $c(n)=0$ for odd $n$. In terms of
geometric sequences,
\[
c=\tfrac12\bigl(E_1+E_{-1}\bigr),\qquad d=E_1-E_0,
\]
because $\frac12(1^n+(-1)^n)$ is $1$ for even $n$ and $0$ for odd $n$. (In $\F_5$ we can divide by $2$: $\frac12=3$,
since $2\cdot3=6\equiv1$.) We look for the solution as a combination of two geometric sequences, and try
\[
b=2E_2-2E_{-1}.
\]
Using $E_s\star E_t=E_{s+t}$,
\[
b\star c=2\cdot\tfrac12\bigl(E_{2+1}+E_{2-1}\bigr)-2\cdot\tfrac12\bigl(E_{-1+1}+E_{-1-1}\bigr)
=E_3+E_1-E_0-E_{-2}.
\]
Since $-2\equiv3\pmod5$, the sequences $E_{-2}$ and $E_3$ are equal, so $b\star c=E_1-E_0=d$. By the uniqueness in
Lemma~\ref{lem:rec}, applied over $\F_5$, we get $a\equiv b$, that is,
\[
a(n)\equiv2\cdot2^n-2\cdot(-1)^n\equiv2\cdot2^n+3\cdot4^n\pmod5,
\]
since $-2\equiv3$ and $-1\equiv4$. Because $2^4=16\equiv1$ and $4^4\equiv1\pmod5$, the remainder depends only on $n$
modulo $4$. Evaluating at $n=0,1,2,3$ gives $2+3=5$, $4+12=16$, $8+48=56$ and $16+192=208$, whose remainders modulo $5$ are
$0,1,1,3$. Only the first is $0$, so $5$ divides $a(n)$ exactly when $4$ divides $n$.
\end{proof}

\section{Checking the results}\label{sec:checks}

A mathematical proof can contain mistakes, so the results were checked in two independent ways.

\para{Numerical check.}
We expanded $(e^x-1)/\cos2x$ as a power series with exact rational arithmetic. This reproduces all 22 terms printed
in the OEIS entry. All $a(n)$ with $n<400$ are integers and satisfy both parts of Theorem~\ref{thm:main}.
The computation was done with the PARI/GP computer algebra system, using code that was written separately from the
formal proofs (file \texttt{verification/independent\_checks.gp} in the accompanying repository~\cite{repo}).
A computation can only test finitely many cases, so it is a check on the definitions and on the indexing, not a proof.

\para{Formal proof in Lean.}
Lean~\cite{lean} is a \emph{proof assistant}: a program in which mathematical statements and their proofs are written
in a precise formal language, and every step is checked mechanically down to the basic rules of logic. All results of
this paper are proved in Lean~4 (version 4.19.0) on top of Mathlib~\cite{mathlib}, Lean's library of formalized
mathematics. The statements are collected in the file \texttt{lean/Entries/A380056.lean} of the repository,
in the namespace \texttt{OEIS.A380056}. There the series $e^x$ and $\cos x$ are Mathlib's formal exponential and cosine series over $\Q$, with $x$ replaced by $2x$ in the cosine. The integer sequence is defined by the recurrence of Lemma~\ref{lem:rec}, and the proofs are in the file \texttt{lean/Families/ExpCosine.lean}.
The table lists where each result of this paper is proved.
\begin{center}\small\renewcommand{\arraystretch}{1.25}
\begin{tabular}{>{\raggedright\arraybackslash}p{0.52\linewidth}>{\raggedright\arraybackslash}p{0.4\linewidth}}
\toprule
Statement in this paper & Lean theorem(s)\\
\midrule
$a(n)=n!\,\coeff{n}{A}$, where $A\cdot\cos2x=e^x-1$ over $\Q$ and $A$ is the only such series; the $a(n)$ are integers & \texttt{integer\_\allowbreak{}egf\_\allowbreak{}definition}, \texttt{source\_\allowbreak{}egf}, \texttt{literal\_\allowbreak{}source\_\allowbreak{}equation}, \texttt{literal\_\allowbreak{}source\_\allowbreak{}unique}, \texttt{integer\_\allowbreak{}coefficients} \\
Corollary~\ref{cor:conj}: the two conjectures & \texttt{first\_\allowbreak{}source\_\allowbreak{}conjecture}, \texttt{second\_\allowbreak{}source\_\allowbreak{}conjecture} \\
Theorem~\ref{thm:main}(b) and its consequences & \texttt{mod\_\allowbreak{}five\_\allowbreak{}exponential\_\allowbreak{}formula}, \texttt{exact\_\allowbreak{}mod\_\allowbreak{}five\_\allowbreak{}pattern}, \texttt{exact\_\allowbreak{}divisibility\_\allowbreak{}by\_\allowbreak{}five}, \texttt{mod\_\allowbreak{}five\_\allowbreak{}period} \\
\bottomrule
\end{tabular}
\end{center}
The Lean project was rebuilt from source for this paper, and Lean was asked to list the basic assumptions (axioms)
that each of these theorems ultimately depends on. Every one depends only on the three standard axioms of Lean's logic
(\texttt{propext}, \texttt{Classical.choice}, \texttt{Quot.sound}); no step is left unproved. A formal proof certifies
exactly the statement that Lean checks. Whether that statement says what the OEIS entry means is a separate question,
so the Lean definitions were also compared, side by side, with the definition in the OEIS entry
(Section~\ref{sec:sequence}).

\section{Discussion}\label{sec:discussion}

\para{Scope.} Theorem~\ref{thm:main} describes the remainders of every $a(n)$ modulo $4$ and modulo $5$, so it settles
both conjectures completely.

\para{Other moduli.} A computer calculation for $n<400$ suggests that the remainders of $a(n)$ are eventually periodic
for other moduli as well:
\begin{itemize}
\item modulo $3$, with period $8$ from $n=1$ on;
\item modulo $7$, with period $48$ from $n=1$ on;
\item modulo $8$, with period $4$ from $n=3$ on (the remainders are $5,1,1,5$ for $n\equiv3,0,1,2\pmod 4$);
\item modulo $16$, with period $8$ from $n=3$ on.
\end{itemize}
We have not proved these observations. The proof of part~(b) suggests where the periods $8=3^2-1$ and $48=7^2-1$ may
come from. Modulo $5$ the number $-4$ is a square ($-4\equiv1=1^2$), which is what made $\cos2x$ a combination of
exponentials. Modulo $3$ and modulo $7$ it is not a square, and a similar argument would need a larger ring in
which $-4$ has a square root.

\section*{About this paper}\label{sec:about}
\para{How the work was done.}
This paper is one of 26 papers on conjectures from the OEIS that came out of an autonomous AI research effort directed
by Dev~D.~Goyal~\cite{repo}. The roles were as follows.
\begin{itemize}
\item \emph{Dev~D.~Goyal} started the project, chose which problems to pursue, asked for proofs, formal verification and
revisions, and decided what to publish. Dev~D.~Goyal did not supply proof strategies or intermediate arguments.
\item \emph{ChatGPT} (OpenAI), using the GPT-6 model with Pro reasoning, worked autonomously from those prompts. It
searched the OEIS for suitable conjectures, found the proofs, wrote the Lean~4 formalization, and wrote a first draft
of this paper.
\item \emph{Claude Opus~5.5}, used with Claude Code, checked the work independently and then rewrote this paper. It
rebuilt the Lean project from source and confirmed that the theorems listed in Section~\ref{sec:checks} rely only on
Lean's standard axioms; compared the Lean definitions with the OEIS definition; recomputed the sequence with separately
written PARI/GP code, compared it with the terms published in the OEIS, and tested every claim of the paper on many
terms; and read every step of the proofs. It then rewrote the text to make it self-contained and readable. The
theorems and the proof strategies are ChatGPT's. In rewriting, Claude Opus~5.5 expanded the explanations and, in
places, wrote out an argument in more detail or by a more elementary route. The Discussion section was added by Claude Opus~5.5;
the further observations in it were checked by computer and are not part of the Lean formalization.
\end{itemize}
Dev~D.~Goyal has verified the results and is responsible for this paper.

\para{Review status.} This paper has not been peer reviewed, and it has not been submitted to a journal or to the OEIS.

\para{Priority.} An OEIS entry that labels a statement as a conjecture shows only that its contributor did not know
a proof; it does not show that no proof existed. During the research only limited searches for earlier work were
made. No claim of priority is made.

\para{Corrections.} Please report any error in the issue tracker of the repository~\cite{repo}, where the Lean sources and the numerical checks are also
available.

\section*{Acknowledgements}
The sequence and both conjectures are due to Paul~D.~Hanna~\cite{oeis}. The asymptotic formula mentioned in the
introduction was added to the entry by Vaclav~Kotesovec.

\begin{thebibliography}{99}\small
\raggedright
\bibitem{oeis} P.~D.~Hanna, entry A380056 (E.g.f.\ $(e^x-1)/\cos(2x)$), \emph{The On-Line Encyclopedia of Integer Sequences}, January 2025. \url{https://oeis.org/A380056}. Text as of October 2026.
\bibitem{oeissec} N.~J.~A.~Sloane, entry A000364 (Euler or secant numbers), \emph{The On-Line Encyclopedia of Integer Sequences}. \url{https://oeis.org/A000364}.
\bibitem{oeisdb} N.~J.~A.~Sloane and The OEIS Foundation Inc., \emph{The On-Line Encyclopedia of Integer Sequences}. \url{https://oeis.org}.
\bibitem{stanley} R.~P.~Stanley, \emph{Enumerative Combinatorics}, Volume~2, Cambridge University Press, 1999. (Chapter~5 treats formal power series, composition and compositional inverses.)
\bibitem{niven} I.~Niven, Formal power series, \emph{American Mathematical Monthly} 76 (1969), 871--889.
\bibitem{lean} L.~de~Moura and S.~Ullrich, The Lean~4 theorem prover and programming language, in \emph{Automated Deduction -- CADE 28}, Lecture Notes in Computer Science 12699, Springer, 2021, 625--635.
\bibitem{mathlib} The mathlib Community, The Lean mathematical library, in \emph{Proceedings of the 9th ACM SIGPLAN International Conference on Certified Programs and Proofs (CPP 2020)}, 367--381. Version used: commit \texttt{c44e0c8ee63c}.
\bibitem{repo} D.~D.~Goyal, \emph{Proofs and corrections of OEIS conjectures, with Lean~4 formalizations}, repository, 2026. \url{https://github.com/dvgyl/oeis-a396843-proof}.
\end{thebibliography}
\end{document}
